\section{Related work}
\label{sec:related}

\subsection{Ethics and governance of AI principles}
\label{sec:related-ethics}

\textbf{From principles to practice.} AI ethics has produced many principles \citep{Jobin2019}. \citet{Mittelstadt2019} argues that principles alone are not enough to ensure ethical AI: in contrast to, e.g., medicine, AI development lacks the shared aims, professional norms and accountability mechanisms that allow principles to guide behaviour effectively. A recent form of AI principles are AI constitutions and model specifications, which are intended to be given to and internalised by an AI model during training \citep{bai2022constitutional,Anthropic2026,OpenAI2025}. Existing work on these documents has focused on which principles should be included \citep{kundu2023specific} and how they might be drafted with public participation \citep{huang2024collective}. Less attention has been paid to what happens when the AI itself is asked to suggest revisions to the document intended to govern its behaviour, which we study empirically in this paper.

\textbf{Human oversight.} Human oversight is both an ethical principle \citep{HLEG2019} and a legal requirement. Article 14 of the EU AI Act \citep{EUAIAct2024} requires that high-risk AI systems be designed so that natural persons can effectively oversee them, including by stopping them. Meaningful human control over autonomous systems is an active area of research. Philosophical accounts hold that control must track the reasons relevant to humans and be traceable to human agents \citep{SantoniDeSio2018}. \citet{Sterz2024} argue that for effective oversight the overseer must (i) have epistemic access to the system, (ii) have self-control and fitting intentions, and (iii) have causal power over the system. Policies requiring human oversight have been criticised for assigning people tasks they are not well suited for \citep{Green2022}. Such accounts focus on what the overseer must fulfil in order to exercise oversight effectively.

From a technical perspective, work on corrigibility asks whether the AI system itself will leave oversight intact \citep{soares2015corrigibility}. For example, it has been shown that it is not sufficient for the human to want to switch off the system: the system must also have an incentive to allow itself to be switched off, which can arise from uncertainty about its own objective \citep{HadfieldMenell2017}.

Our work on written oversight commitments is situated between these two areas of research. It deals with oversight commitments that the developer writes into the system's own specification, and with the question of whether these commitments survive revision by the system itself.

\textbf{Political neutrality.} It has been shown that language models exhibit political opinions that are unevenly aligned with the opinions of various groups of people \citep{Santurkar2023}, and that political biases learned in pretraining transfer to downstream tasks \citep{Feng2023}. It is therefore an open question whether and how models ought to express political opinions. Our constitution includes a restriction against doing so, so the stability of that restriction under reflection has a direct bearing on this question.

\textbf{Constitutions and principles.}
Constitutional AI trains models against a written list of principles, using model-generated critiques and revisions \citep{bai2022constitutional}. Later work looks at the choice of principles \citep{kundu2023specific}, draws them from the public \citep{huang2024collective}, and audits how closely models follow them \citep{jakkli2026constitutions}. Deliberative alignment trains models to reason explicitly over a written safety specification \citep{guan2024deliberative}. All this work uses principles that humans wrote. Instead, we let the governed model itself rewrite its own principles iteratively, and measure what survives.

\textbf{Iterated transmission and self-refinement.}
Repeated transmission of text through LLMs results in systematic drift towards LLM-specific attractors \citep{perez2024telephone,mohamed2025broken,geng2026markovian}. Self-refinement applied recursively also leads to convergence to fixed points preferred by the model, with most of the change occurring in the first few rounds, in repeated refinement of scientific abstracts \citep{wu2026converge}. We utilize this convergence as a measurement tool. Paraphrase and neutral-edit conditions, along with matched content categories, allow us to disentangle value-specific erosion from general drift of text.

\textbf{Self-evolving agents and value drift.}
Agents that rewrite their own controllers, memories or skills under task pressure show eroded safety behaviour \citep{das2026seabench}. Reward-driven self-evolution can tip agents away from alignment \citep{han2025alignmenttipping}. Self-evolving agents can also drift through changes to their memory, tools and workflows \citep{shao2025misevolve}. Misaligned goals written into persistent memory propagate to later agents \citep{das2026selfpropagating}. Goal drift has also been documented over long agentic episodes \citep{arike2025goaldraft} and across inherited contexts \citep{menon2026inherited}. Our setting removes task pressure, reward and memory, isolating reflection on stated values.

\textbf{Oversight evaluations and measurement.}
Our analysis also builds on behavioural evaluations of corrigibility, including shutdown resistance \citep{schlatter2025shutdown}, alignment faking \citep{greenblatt2024alignmentfaking}, in-context scheming \citep{meinke2024scheming}, and corrigibility failures in computer-use agents \citep{tien2026rogue}. Methodologically, our work borrows from discrete-time survival analysis, recently applied to assess multi-turn consistency in language-model agents \citep{bogdanov2026survival}, and from concerns about the reliability and validity of LLM judges \citep{norman2026reliability,mukherjee2026geometry,zheng2023judging}, which motivated our two-judge approach.
